% =========================================================================
% Paper 1: CASPER — Controlled Assessment of Strategic Preference
% Elicitation in Conversational Recommendation
%
% Long journal format (thesis-includable). Draft v0.1 — results sections
% contain placeholders marked [RESULTS]. Compile: pdflatex + bibtex
% (references.bib in this folder, merged from thesis bibliography).
% =========================================================================
\documentclass[11pt]{article}

\usepackage[margin=1.1in]{geometry}
\usepackage{amsmath, amssymb}
\usepackage{booktabs}
\usepackage{multirow}
\usepackage{graphicx}
\usepackage[hidelinks]{hyperref}
\usepackage[comma,authoryear]{natbib}
\usepackage{xcolor}
\newcommand{\todo}[1]{\textcolor{red}{[TODO: #1]}}
\newcommand{\results}[1]{\textcolor{blue}{[RESULTS: #1]}}

\title{CASPER: Controlled Assessment of Strategic Preference Elicitation
in Conversational Recommendation}

\author{Aleksandra Makarova \and Muhammad Shahzad \and Xia Hong \and Martin Lester\\
\small Department of Computer Science, University of Reading, UK}

\date{Draft v0.1 --- June 2026}

\begin{document}
\maketitle

\begin{abstract}
Conversational recommender systems (CRS) are increasingly built around large
language models (LLMs), yet their defining capability --- strategically
eliciting user preferences --- is asserted far more often than it is
measured. Published evaluations embed elicitation inside whole-system
comparisons, where its effect is confounded with recommendation and
generation quality, and the heterogeneous approaches now in play
(graph-pruned Q-networks, Bayesian acquisition, prompted LLMs, learned
policies) have never been compared on equal terms. We introduce CASPER, a
controlled testbed that isolates elicitation strategy: a fixed,
polarity-aware recommendation model serves as the measurement instrument,
simulated users answer verifiably from held-out MovieLens ground truth, and
policies are scored on recommendation-relevant information gained per
conversational turn. We first show that instrument design is decisive.
Under common configurations, the instrument itself fails in ways that
silently invalidate measurement: models trained on full preference profiles
treat the few-reveals regime of an actual dialogue as out-of-distribution,
and sentence-embedding-based user representations conflate preference with
anti-preference --- in our audit, a widely-used configuration ranked
franchise sequels \emph{higher} when the user disliked the franchise anchor.
We propose acceptance tests that any candidate instrument must pass
(polarity direction, reveal monotonicity, signal-to-noise), and a
reveal-subset training augmentation that passes them. We then benchmark
eight elicitation policies spanning random and popularity baselines, greedy
information gain, prompt-only LLM questioning in three styles, and a
reinforcement-learned bot-play policy. \results{headline finding, e.g.\
``Prompt-only LLM questioners underperform greedy information gain by X\%
AUAC at equal turn budgets''}. The testbed protocol, simulator, instrument
acceptance suite, and all reference policies are released to enable
measurement-grounded progress on strategic elicitation.
\end{abstract}

\section{Introduction}
\label{sec:intro}

The promise of a conversational recommender system (CRS) over a traditional
one is interactivity: rather than inferring preferences solely from logged
behaviour, the system can \emph{ask}. What it asks, and when, determines how
much of the user's limited conversational patience is converted into
recommendation quality. This capability --- \emph{strategic preference
elicitation} --- has been the explicit objective of a decade of research,
from early reinforcement-learning dialogue managers \citep{sun2018conversational,
lei2020estimation} through graph-constrained policy learning
\citep{lei2020interactive, deng2021unified} to, most recently, large language
models (LLMs) prompted to conduct preference interviews
\citep{li2023eliciting, he2023large}.

Yet a researcher who asks the simple question \emph{``which of these
approaches asks better questions?''} will find no answer in the literature.
Published evaluations measure whole systems: a dialogue agent's questions,
its recommendation model, its language understanding, and its generation
quality are exercised together, and a single end metric --- success rate,
hit ratio, NDCG at dialogue end --- absorbs all of them
\citep{jannach2021survey, wang2023rethinking}. A system may score well
because it asks incisive questions, or because its recommender extracts more
from the same answers, or because its utterances are parsed more reliably;
the published numbers cannot tell these apart. The confound has become more
acute, not less, as the field has diversified: the candidate elicitation
mechanisms now include discrete attribute-selection policies trained with
deep Q-networks \citep{deng2021unified, zhao2023hierarchical}, single-step
Bayesian acquisition over item utilities \citep{austin2024bayesian},
unoptimised free-form LLM interviewing \citep{li2023eliciting}, and hybrid
architectures in which a learned planner selects strategies that a frozen
LLM verbalises \citep{deng2023plug, zhao2025reinforced}. These mechanisms
share no common action space, no common training signal, and no common
evaluation protocol.

The absence of controlled measurement has a second, subtler consequence:
the field's motivating claims rest on weak evidence. The widely-repeated
assertion that prompt-only LLMs are poor elicitation strategists ---
recommending prematurely, asking unfocused questions --- originates in
qualitative observations and in motivation sections of papers proposing
their own fixes \citep{he2023large, zhao2025reinforced, kostric2025should}.
We are aware of no controlled quantification of how strategically LLMs
actually question, relative to principled non-learned baselines such as
greedy information gain. If the gap is large, learned elicitation policies
have a clear target; if it is small, a substantial research agenda needs
re-aiming. Either way, the field should know.

This paper introduces \textbf{CASPER} (\emph{Controlled Assessment of
Strategic Preference Elicitation in Recommendation}), a testbed that
isolates elicitation strategy as the single experimental variable. CASPER
rests on three design commitments:

\begin{enumerate}
  \item \textbf{A fixed, audited measurement instrument.} A single
  pre-trained recommendation model maps revealed preferences to predicted
  ratings for all policies alike. Crucially, we treat the instrument itself
  as an object of audit: we define \emph{acceptance tests} --- polarity
  direction, reveal monotonicity, and signal-to-noise requirements --- that
  an instrument must pass before any policy comparison performed with it
  can be meaningful. We show that plausible, previously-used instruments
  fail these tests in instructive ways (Section~\ref{sec:instrument}).
  \item \textbf{A verifiable simulated user.} Simulated users answer
  entity-level questions deterministically from held-out ground-truth
  profiles. Unlike free-form LLM simulators \citep{wang2023rethinking},
  every simulated statement is auditable against the profile, which is
  essential when the measured quantity is information gain: a simulator
  that hallucinates a preference injects unmeasurable noise into precisely
  the signal under study.
  \item \textbf{A shared discrete action interface.} Every policy ---
  random, information-theoretic, prompted-LLM, or learned --- selects each
  question from the same fixed slate of entities (movies, genres, actors,
  directors). Strategy is thereby the only varying factor; differences in
  action vocabulary, language generation, and parsing are excluded by
  construction.
\end{enumerate}

Within this testbed we benchmark eight reference policies: uniform random;
population-popularity ordering; greedy expected information gain (one-step
entropy reduction under the instrument, the strongest non-learned baseline);
three prompt-only LLM questioners of increasing sophistication (a vanilla
asker, a chain-of-thought strategist, and a GATE-style interviewer
\citep{li2023eliciting}); and a reinforcement-learned bot-play policy
following our earlier published method \citep{makarova2024learning}.

\paragraph{Contributions.}
(1)~A formal statement of the elicitation measurement problem and a
controlled protocol that solves it (Section~\ref{sec:problem},
\ref{sec:testbed}).
(2)~An instrument acceptance suite, together with an empirical audit showing
that two plausible instrument configurations fail it --- one because partial
reveals are out-of-distribution, one because sentence-embedding user
representations break preference polarity --- and a reveal-subset training
augmentation that repairs the failure (Section~\ref{sec:instrument}).
(3)~The first controlled, like-for-like comparison of LLM questioning
against information-theoretic and learned elicitation baselines
(Section~\ref{sec:results}).
(4)~A released implementation comprising the testbed, simulator, acceptance
suite, reference policies, and evaluation scripts.
\todo{add repository URL on release}

\section{Related Work}
\label{sec:related}

\subsection{Learned Question Selection in CRS}

Reinforcement-learning approaches to conversational recommendation have
treated question selection as a sequential decision problem over discrete
action spaces since \citet{sun2018conversational}, whose policy chooses
among a fixed set of item facets, and \citet{lei2020estimation}, whose
Estimation--Action--Reflection cycle alternates attribute questions with
recommendations in a multi-round protocol. Subsequent systems concentrated
on taming the size of the discrete space: SCPR \citep{lei2020interactive}
restricts candidate attributes to graph neighbours of confirmed preferences
and reduces the learned decision to a binary ask-or-recommend choice;
UNICORN \citep{deng2021unified} scores a pruned candidate set of items and
attributes with a dueling Q-network; hierarchical refinements separate
when-to-ask from what-to-ask \citep{zhao2023hierarchical} or generate
chains of attribute choices \citep{fan2024chain}. Throughout this lineage,
evaluation is conducted against rule-based simulators with template
questions, and --- material for the present paper --- the recommendation
model, simulator, and policy are co-designed and co-reported, so the
elicitation contribution is never measured in isolation.

\subsection{LLMs as Conversational Recommenders and Elicitors}

Zero-shot LLMs outperform fine-tuned CRS models on established benchmarks
\citep{he2023large}, with the advantage concentrated in content and context
knowledge rather than collaborative signals, and characteristic failure
modes including popularity bias and out-of-catalogue recommendation.
\citet{wang2023rethinking} showed that static-replay benchmarks
systematically penalise interactive behaviour and introduced LLM-based user
simulation (iEvaLM), which has become the dominant interactive evaluation
protocol --- at the cost of making the simulator's biases, catalogue
knowledge, and cooperativeness part of the experimental design. On the
elicitation side specifically, GATE \citep{li2023eliciting} demonstrated
that prompted LLMs can conduct free-form preference interviews that beat
user-written specifications; PEBOL \citep{austin2024bayesian} added
decision-theoretic structure, maintaining Bayesian beliefs over item
utilities with the LLM as a natural-language acquisition function, and
showed that this beats monolithic prompted elicitation --- evidence that
explicit strategy helps even atop strong language models. Hybrid
architectures train a compact policy to select dialogue strategies that a
frozen LLM verbalises \citep{deng2023plug}, most recently for CRS with a
thirteen-strategy inventory including three elicitation types
\citep{zhao2025reinforced}. \citet{zhang2025modeling} optimise clarifying
questions at the token level via preference labels derived from simulated
future turns.

None of these works measures elicitation strategy under a fixed instrument
and verifiable simulator; the comparisons reported are whole-system, and
the LLM-elicitation-weakness claim that motivates several of them remains,
to our knowledge, unquantified. CASPER is designed to supply exactly this
missing measurement layer; it is complementary to, not competitive with,
the systems above, each of which could be entered as a policy.

\subsection{Preference Elicitation Beyond CRS}

Preference elicitation predates CRS in decision theory, where questions are
selected to reduce uncertainty about a utility function by minimax regret
\citep{viappiani2009regret} or Bayesian updating \citep{mangili2020bayesian},
and in interactive recommendation via bandit formulations over absolute and
comparative item questions \citep{christakopoulou2016towards}. Our greedy
information-gain baseline transplants this tradition's one-step acquisition
principle onto the testbed's instrument, providing the natural
non-learned reference point for all learned and prompted policies.

\subsection{Evaluation Methodology}

Simulation-based CRS evaluation spans rule-based attribute-answering
simulators \citep{lei2020estimation, deng2021unified}, crowd-sourced
proxies \citep{sun2018conversational, kang2019recommendation}, and LLM
role-play \citep{wang2023rethinking}. Critiques of CRS evaluation
\citep{jannach2022evaluating} emphasise weak correlation between automatic
metrics and user-perceived quality, and the difficulty of comparing
multi-component systems. Our contribution to this conversation is narrower
and sharper: for the specific quantity \emph{information gained per
question}, we show the measurement instrument's statistical properties
(polarity sensitivity, reveal-regime coverage, signal-to-noise) determine
whether any comparison is meaningful at all, and we provide the audit
machinery to establish those properties before benchmarking.

\section{The Measurement Problem}
\label{sec:problem}

\subsection{Strategic Elicitation as a Sequential Decision Problem}

At each turn $t$ of an elicitation dialogue, the system holds a record
$R_t = \{(e_i, a_i)\}_{i<t}$ of entities asked about and the user's answers,
selects a next entity $e_t \in \mathcal{E} \setminus \{e_i\}_{i<t}$ from a
fixed slate $\mathcal{E}$, and receives an answer
$a_t \in \{\textsf{liked}, \textsf{disliked}, \textsf{unknown}\}$. A fixed
recommendation model $\phi$ --- the \emph{instrument} --- maps the record to
predicted ratings $\hat{y} = \phi(R_t)$ over a movie catalogue
$\mathcal{M}$, and the quality of elicitation so far is the agreement
between $\hat{y}$ and the user's held-out true ratings $y^*$:
\begin{equation}
\mathrm{Acc}_t \;=\; \frac{1}{|\mathcal{M}_u|} \sum_{m \in \mathcal{M}_u}
\mathbb{1}\!\left[\, (\hat{y}_m > 0.5) = y^*_m \,\right],
\label{eq:acc}
\end{equation}
where $\mathcal{M}_u \subseteq \mathcal{M}$ are the movies the user has
actually rated. A \emph{policy} $\pi$ maps the dialogue record to the next
entity. We compare policies by their accuracy trajectories
$(\mathrm{Acc}_0, \ldots, \mathrm{Acc}_T)$: the per-turn curve, its final
value, and its mean over turns (area under the accuracy curve, AUAC), with
$\mathrm{Acc}_0$ the instrument's prior before any question.

This formulation makes the isolation requirements explicit. The instrument
$\phi$ and the slate $\mathcal{E}$ are shared by all policies; the answer
function is a deterministic lookup in the user's ground-truth profile; the
only free variable is $\pi$. Anything a policy might gain through a better
recommender, a richer action vocabulary, or more reliable language
processing is excluded by construction.

\subsection{Why Whole-System Evaluation Cannot Answer the Question}

Consider the standard multi-round CRS protocol \citep{lei2020estimation}:
the system asks or recommends, the (simulated) user responds, and success
rate at a turn budget is reported. Three quantities vary together: the
policy's questions, the recommender that converts answers into rankings,
and the simulator's response behaviour. Improvements in any one move the
end metric. The literature's standard remedy --- ablations within a single
system --- establishes that a component matters to \emph{that} system, not
how competing elicitation mechanisms compare across systems. The problem
compounds for LLM pipelines, where question quality, parsing reliability,
and recommendation knowledge are entangled within a single model's
behaviour, and for LLM-simulated users, whose answering behaviour is itself
policy-dependent (a fluent questioner may elicit longer, more revealing
responses from the simulator than a templated one --- a property of the
simulator, not of the user population the simulation stands in for).

\subsection{Why the Instrument's Statistics Decide Measurability}
\label{sec:snr_theory}

Beyond isolation, measurement requires the instrument to be \emph{sensitive
to the signal being measured}. Two properties are necessary:

\paragraph{Polarity.} The instrument must respond in opposite directions to
\textsf{liked} and \textsf{disliked} answers about the same entity.
An instrument that conflates them cannot reward a policy for eliciting
disambiguating information, and policies that ask polarity-resolving
questions will measure no better than ones that do not. This requirement is
not hypothetical: representations that encode preferences as text in a
sentence-embedding space assign high similarity to ``likes $X$'' and
``dislikes $X$'' --- in our preliminary system, cosine similarities of
0.75--0.78 --- and recommenders built on them produce largely overlapping
recommendations for opposite preference profiles
(Section~\ref{sec:instrument_audit}).

\paragraph{Signal-to-noise.} The per-question accuracy gain must be
resolvable against between-user variance at the achievable sample size. In
preliminary work with a ranking-based reward over small candidate pools, we
measured per-turn metric deltas of $\sim$0.005 against a per-user standard
deviation of $\sim$0.20 --- a per-turn signal-to-noise ratio of
$\sim$0.025, at which neither reinforcement learning nor policy comparison
is statistically meaningful at realistic budgets. The same arithmetic
governs evaluation: with $n$ users, the standard error of a mean accuracy
difference scales as $\sigma\sqrt{2/n}$; an instrument with per-reveal
gain $g$ and user-level dispersion $\sigma$ needs
$n \gtrsim 2(z\sigma / kg)^2$ users to resolve a $k$-question quality
difference at confidence $z$. We therefore report, as part of instrument
acceptance, the empirical gain and dispersion that determine the user
sample needed.

\section{A Case Study in Measurement Failure}
\label{sec:casestudy}

Before describing the testbed, we document the experience that motivated
it: a reinforcement-learning research programme on elicitation that stalled
for reasons that were ultimately about measurement, not learning. We report
it both as evidence for the claims of Section~\ref{sec:problem} and because
we believe the pattern is common but rarely published.

\subsection{Three Failed Training Campaigns}

Building on \citet{makarova2024learning}, we attempted to train an
elicitation policy with a continuous semantic action space, rewarded by
per-turn improvements in NDCG@10 of an LLM-based recommender over a
candidate pool, against an LLM-simulated user. Three campaigns failed in
sequence. The first produced \emph{negative} mean rewards
($-0.25$ over the final ten episodes): the evaluation pool was so saturated
with popular ground-truth items that the recommender's zero-knowledge
baseline was already high, and revealing genuine preferences frequently
\emph{lowered} the metric. The second collapsed into repetition once the
pool configuration was corrected. The third, with a within-episode
similarity penalty, converged to a fixed question sequence delivered
identically to every user: a single concept accounted for 40\% of all
second-turn questions, and only 35\% of 783 question selections were
unique across 90 episodes. At no point did any campaign learn
user-adaptive questioning.

\subsection{Quantifying the Reward Channel}

Post-hoc diagnostics located the failure in the measurement channel rather
than the policy. Table~\ref{tab:pool_baseline} shows the zero-knowledge
baseline of the NDCG-based reward as a function of two protocol parameters:
candidate pool size and the minimum number of ground-truth items required
per evaluation user. The configuration used in training (pool 100, minimum
25 targets) yields a baseline of 0.52 --- with $\sim$39\% of the pool being
ground truth, popularity guessing scores well before any question is asked
--- leaving almost no headroom for elicitation to be rewarded.

\begin{table}[t]
\centering
\caption{Zero-knowledge NDCG@10 baseline by evaluation protocol parameters
(LLM recommender, $n = 50$ users per cell, per-user std $\approx 0.20$).
Protocol choices, not elicitation quality, dominate the metric.}
\label{tab:pool_baseline}
\small
\begin{tabular}{lccccc}
\toprule
Pool size & \multicolumn{5}{c}{Minimum ground-truth targets per user} \\
 & 5 & 10 & 15 & 20 & 25 \\
\midrule
100 & 0.315 & 0.412 & 0.460 & 0.508 & \textbf{0.519} \\
200 & 0.275 & 0.295 & 0.420 & 0.430 & 0.463 \\
500 & 0.253 & 0.286 & 0.295 & 0.377 & 0.417 \\
\bottomrule
\end{tabular}
\end{table}

\begin{table}[t]
\centering
\caption{Effect of revealed-preference type on the same reward channel
(pool 100, 20 reveal steps, identical $n = 30$ users per condition;
95\% CI $\approx \pm 0.074$). Only the title-only decline and the mixed
improvement exceed noise.}
\label{tab:pref_type}
\small
\begin{tabular}{lcccc}
\toprule
Revealed preferences & Baseline & Peak (step) & Final & $\Delta$ \\
\midrule
Movie titles only & 0.323 & 0.394 (3) & 0.215 & $-0.109$ \\
Attributes only & 0.323 & 0.429 (8) & 0.348 & $+0.025$ \\
Mixed (interleaved) & 0.323 & 0.451 (5) & 0.425 & $+0.102$ \\
\bottomrule
\end{tabular}
\end{table}

Table~\ref{tab:pref_type} probes the channel's response to the
\emph{content} of revealed preferences under a scripted (policy-free)
reveal schedule. The response is non-monotonic and type-dependent ---
revealing movie titles alone eventually drives the metric \emph{below} its
zero-knowledge baseline --- and the per-step deltas ($\sim$0.005 on
average) sit against a per-user standard deviation of $\sim$0.20. At that
signal-to-noise ratio ($\sim$0.025 per step), the per-turn reward an
individual training episode delivers is statistically indistinguishable
from noise, and three campaigns of policy learning amounted to fitting
that noise. No property of the policy class, exploration schedule, or
reward shaping --- the knobs we had been turning --- could have repaired
this.

\subsection{Lessons}

The case study yields the design requirements of
Section~\ref{sec:testbed}: (i) the measurement instrument must be audited
\emph{before} policy work, against explicit acceptance criteria; (ii) the
simulated user must be verifiable, so that reward noise cannot enter
through hallucinated or inconsistent answers; (iii) the per-turn signal
and its dispersion must be measured and reported, because they determine
both learnability and the user sample needed for evaluation; and (iv) the
protocol parameters that Table~\ref{tab:pool_baseline} shows to dominate
the metric must be fixed and shared across all compared policies.

\section{The CASPER Testbed}
\label{sec:testbed}

\subsection{Data and Entity Slate}

CASPER instantiates the protocol on MovieLens-25M \citep{harper2015movielens}.
The entity slate $\mathcal{E}$ comprises 187 items: the 100 most-rated
movies, 17 genres, the 50 most-frequent lead actors and 20 directors of
those movies (cast and crew from TMDB). Movie ground truth is binarised at
the conventional threshold (rating $\geq 4$ is \textsf{liked}); attribute
ground truth is derived per user as the mean rating of the user's rated
movies matching the attribute, binarised at the same threshold; entities
without support in the user's history answer \textsf{unknown}. We restrict
to users with at least 25 ratings on the slate movies (64{,}639 users), so
that accuracy (Eq.~\ref{eq:acc}) is estimated over a meaningful per-user
movie set.

\paragraph{User partitions.} A fixed shuffle (seed 42) divides users 80/20.
The 80\% partition trains instruments and learned policies; the first
1{,}500 users of the 20\% partition serve as the instrument's
model-selection set; the remaining $\sim$11{,}400 users are reserved
exclusively for testbed evaluation and touch no training or selection step
of any component.

\subsection{The Verifiable Simulated User}

A simulated user is a ground-truth profile vector over the slate. Asked
about entity $e$, it answers \textsf{liked} or \textsf{disliked} by lookup,
or \textsf{unknown} where the profile has no support. The simulator is
deterministic, free, and auditable: every answer can be checked against the
profile, and the information content of a dialogue is exactly the set of
(entity, polarity) pairs revealed. We deliberately trade away linguistic
naturalism --- the simulator emits labels, not utterances --- because the
measured quantity is information gain, and any generative variability in
answering would be confounded with the policies' questioning behaviour.
Section~\ref{sec:limitations} discusses what this choice does and does not
license.

\subsection{The Measurement Instrument and Its Acceptance Tests}
\label{sec:instrument}

The instrument is an extrapolation recommender in the style of our earlier
work \citep{makarova2024learning}: an LSTM with attention consumes the full
slate, each position carrying the item's identity and a rating one-hot
$[\textsf{disliked}, \textsf{liked}, \textsf{unseen}]$, and predicts rating
logits for every movie; revealed entities carry their answer, all others
\textsf{unseen}. This architecture encodes preference polarity
\emph{structurally} --- as an input channel orthogonal to item identity ---
rather than textually, avoiding by construction the sentence-embedding
conflation discussed above.

\subsubsection{Why a Learned Instrument Rather Than Matrix Factorization}
\label{sec:why_not_mf}

The classical recommender for cold-start elicitation is matrix
factorization (MF) with new-user fold-in: latent item factors are learned
offline, and a new user's answers are converted to a profile by a ridge
least-squares solve against the factors of the answered items
\citep{golbandi2011adaptive, zhou2011functional}. MF is the natural
reference point, and we include it as such (Section~\ref{sec:mf_reference}).
We nonetheless adopt a learned set instrument for four reasons specific to
the quantity this testbed measures --- recommendation-relevant information
gained per \emph{question} --- none of which MF provides:

\begin{enumerate}
  \item \textbf{Answerability is a first-class output.} Elicitation value
  is governed as much by whether a user \emph{can} answer a question as by
  how informative the answer would be (Section~\ref{sec:casestudy}); on
  real catalogues the binding constraint is answer rate, not raw
  informativeness. MF predicts utility, not answer-availability. The
  instrument's second head predicts $P(\textsf{rated}\mid R_t)$ directly,
  so the testbed can both score answerability-aware strategies and
  decompose a policy's gain into ``asks answerable things'' versus ``asks
  informative things'' --- a decomposition unavailable from a utility-only
  recommender.
  \item \textbf{A heterogeneous entity slate in one model.} Questions range
  over movies, genres, actors and directors; MF factorizes a user--item
  matrix and must be extended (side-feature or tensor variants) to score
  non-item entities. The set instrument consumes every entity uniformly as
  an (identity, polarity) token, so attribute and item questions are
  measured on one footing.
  \item \textbf{Open-vocabulary extensibility.} The companion
  continuous-action work asks via free-form, LLM-verbalised probes mapped
  into an embedding space; a recommender indexed by catalogue ID cannot
  score an off-catalogue or generated probe, whereas a set encoder over
  (embedding, polarity) tokens extends to unseen concepts. The instrument
  is thus the same measurement device across both papers.
  \item \textbf{Amortized, differentiable partial-reveal inference.} A
  single forward pass maps any reveal set to a full belief, versus a
  per-user re-solve, and the belief is differentiable --- usable directly
  as the state and reward channel of a learned policy without a
  non-differentiable fold-in step in the loop.
\end{enumerate}

These advantages are only worth having if the learned instrument is at
least as good a \emph{recommender} as MF; otherwise a weak instrument would
suppress every policy's measured gain and the comparison would be
uninformative. We therefore treat parity-with-MF as an acceptance
requirement: the instrument must match an MF reference on standard
recommendation metrics (NDCG@10, Recall@10, MRR) at matched reveal budgets,
\emph{and} elicitation must produce gains under the instrument no smaller
than under MF, before any of the four advantages above are claimed.
Section~\ref{sec:mf_reference} reports this comparison.

We require any candidate instrument to pass four acceptance tests
\emph{before} it is used to compare policies:

\begin{description}
  \item[T1 (attribute polarity).] For each attribute with at least three
  matching slate movies, the mean rank of matching movies must improve when
  the attribute is \textsf{liked} versus \textsf{disliked}. Pass: correct
  direction for $\geq 80\%$ of attributes.
  \item[T2 (reveal monotonicity).] Mean held-out accuracy must increase
  monotonically with the number of revealed preferences (Spearman
  $\rho > 0.9$ across reveal counts $1$--$50$).
  \item[T3 (franchise flip).] For each film-franchise anchor in the slate
  (14 franchises), flipping the anchor's rating from \textsf{liked} to
  \textsf{disliked} must move the franchise's other movies down the
  ranking. Pass: correct direction for $\geq 70\%$ of related movies.
  \item[T4 (signal-to-noise).] The accuracy gain over the first 20 reveals,
  divided by the between-user standard deviation at 10 reveals, is
  reported and determines the evaluation sample size
  (Section~\ref{sec:snr_theory}).
\end{description}

\subsubsection{An Audit of Plausible Instruments}
\label{sec:instrument_audit}

We audited two instrument configurations carried over from our earlier
system, both trained to extrapolate full preference profiles
(Table~\ref{tab:audit}). The first (\textsc{OneHot}) uses learned item
embeddings; the second (\textsc{Concept}) replaces them with frozen
SBERT \citep{reimers2019sentence} text embeddings of the item names.

Both fail, and the failure modes are instructive.
\textsc{OneHot} exhibits genuine polarity --- on T3, disliking a franchise
anchor demotes its sequels by $\sim$20 of 100 rank positions, 91.7\% correct
direction --- and clean monotonicity (65.9\% accuracy at one reveal rising
to 91.6\% at fifty, $\rho = 1.0$). But on T1 it scores only 68.1\%: trained
exclusively on \emph{complete} profiles, the few-reveals regime that an
actual dialogue inhabits is out-of-distribution, and the response to a
single revealed attribute is unreliable. Measured naively, the instrument
would discriminate policies only weakly in exactly the turn range
($\leq$15 questions) the testbed cares about; in a five-turn pilot, even a
greedy information-gain policy moved accuracy by less than 0.011.
\textsc{Concept} fails outright: 51.1\% on T1 (chance) and 16.7\% on T3 ---
\emph{worse} than chance, ranking franchise sequels \emph{higher} when the
anchor is disliked. Frozen text embeddings of item names, projected and
mixed by a sequence encoder trained on full profiles, do not preserve the
polarity channel.

\begin{table}[t]
\centering
\caption{Instrument audit across four generations. \textsc{RS} =
reveal-subset training (Section~\ref{sec:reveal_subset}); \textsc{TR} =
taste-relative attribute labels (Section~\ref{sec:harshness}). DiD =
difference-in-differences attribute polarity (matching movies must move
more than non-matching ones). Each generation fixes the previous failure
and exposes the next. \results{fill final column from v3}}
\label{tab:audit}
\small
\begin{tabular}{lcccc}
\toprule
 & \textsc{OneHot} & \textsc{Concept} & \textsc{+RS} & \textsc{+RS+TR (ours)} \\
\midrule
T1a attr.\ polarity, rank ($\geq$80\%) & 68.1\% & 51.1\% & 51.1\% & \textbf{100\%} \\
T1b attr.\ polarity, DiD ($\geq$80\%) & 72.3\% & 59.6\% & 46.8\% & \textbf{100\%} \\
\quad genres only & 100\% & 56.3\% & 37.5\% & \textbf{100\%} \\
T1c liked/disliked top-10 overlap & 41.2\% & 49.8\% & \textbf{0.7\%} & 26.1\% \\
T2 monotonicity ($\rho$) & 1.00 & 1.00 & --- & \textbf{1.00} \\
T3 franchise flip, rank ($\geq$70\%) & \textbf{91.7\%} & 16.7\% & 83.3\% & 83.3\% \\
\quad DiD form & 62.5\% & 25.0\% & 87.5\% & \textbf{100\%} \\
T4 SNR (gain$_{1\to20}$/$\sigma_{10}$) & 0.90 & 2.72 & --- & 1.11 \\
Accuracy @ 1 / 50 reveals & .659/.916 & .508/.819 & --- & \textbf{.697/.852} \\
Dialogue regime in-distribution & no & no & yes & yes \\
\midrule
Accepted & no & no & no & \textbf{yes} \\
\bottomrule
\end{tabular}
\end{table}

\subsubsection{Reveal-Subset Training}
\label{sec:reveal_subset}

The repair follows from the diagnosis. We retrain the \textsc{OneHot}
architecture with \emph{reveal-subset augmentation}: each training sample
reveals only a random subset of $k$ rated items ($k$ log-uniform between 1
and the profile size; 15\% of samples reveal attributes only), with the
training target remaining the full rated profile. Every input regime the
testbed produces --- from a single revealed genre to a long mixed dialogue
--- is thereby in-distribution. Slate order is shuffled per sample and the
loss is taken at every sequence position, so each prefix functions as a
further, smaller reveal set. Architecture, data, splits, and optimisation
are otherwise unchanged.

\subsubsection{The Rater-Harshness Collapse, and Taste-Relative Labels}
\label{sec:harshness}

Reveal-subset training repaired the regime problem dramatically at the
level of preference \emph{combinations} --- the top-10 recommendations for
a set of attributes all-liked versus all-disliked, which had overlapped by
41\% under the original instrument, became essentially disjoint (0.7\%,
row T1c of Table~\ref{tab:audit}) --- while franchise-anchor flips
remained correct (83\%). But the difference-in-differences test exposed a
new failure: the response to a \emph{single} revealed attribute had
collapsed to chance (46.8\% overall; 37.5\% on genres), even though the
original instrument had scored 100\% on genres by the same test.

Inspection of raw predictions explains the collapse, and the explanation
is a finding in its own right. Conditioning on ``Horror: disliked'', the
retrained instrument lowered its scores for the slate's five horror
classics from $\approx$0.73 to $\approx$0.46 --- the correct direction ---
but lowered \emph{every other movie even further}, so that in rank terms
the horror films \emph{rose} from ranks 39--93 into the top 6. The model
had learned that the derived label ``dislikes Horror'' chiefly identifies
a \emph{harsh rater}: in a slate composed of beloved classics
(\emph{The Silence of the Lambs}, \emph{Alien}, \emph{The Shining}), a
user whose mean rating of these films falls below the liking threshold is,
overwhelmingly, a user who rates everything severely. The signal is real
and genuinely predictive of held-out ratings --- which is exactly why a
well-trained model learns it in preference to the much weaker
attribute-specific taste signal. Under full-profile training this
shortcut is unavailable (dozens of co-revealed movie ratings already pin
down the user's severity), which is why the original instrument had
learned genuine attribute polarity instead.

The repair therefore belongs in the \emph{label construction}, not the
model. We re-derive attribute ground truth \emph{taste-relatively}: an
attribute is \textsf{liked} (\textsf{disliked}) only if the user's mean
rating over matching slate movies exceeds (falls below) \emph{that user's
own overall slate mean} by a margin of 0.25, with at least three matching
rated movies; otherwise the attribute answers \textsf{unknown}. Severity
is thereby removed from the attribute channel by construction, on both
the training side and the simulator side (which derives answers from the
same labels). This also improves the simulator's face validity: a user
who simply rates everything low is no longer described as ``disliking''
every genre. \results{v3 acceptance column in Table~\ref{tab:audit}; if
T1b still $<$80\%, report the auxiliary polarity-contrast loss variant}

The three-generation arc --- full-profile training fails the dialogue
regime; regime-correct training collapses onto rater severity; severity
must be removed from the labels --- illustrates the paper's central
methodological claim in miniature: every one of these failures would have
been invisible in a whole-system evaluation, and each would have silently
distorted any policy comparison run on top of it.

\subsubsection{Matrix-Factorization Reference and Parity Check}
\label{sec:mf_reference}

To establish that the learned instrument is not merely a weaker substitute
for the classical recommender --- which would invalidate every downstream
comparison --- we benchmark it against a matrix-factorization reference
with new-user fold-in, the standard cold-start recommender of
\citet{golbandi2011adaptive, zhou2011functional}. The reference learns
$d$-dimensional item factors and per-item biases on the training partition;
a test user's profile is the ridge least-squares fold-in of the revealed
(item, polarity) answers against those factors, and predictions are
$\hat{y}_m = \bar{y} + b_m + u^\top q_m$. Both recommenders are evaluated on
the identical held-out users and reveal splits, under two lenses:
\emph{(i) recommendation quality}, ranking each held-out liked item against
all non-interacted movies (NDCG@10, Recall@10, MRR) and reconstructing the
held-out preference vector (Brier, ROC-AUC), at full reveal; and
\emph{(ii) elicitation headroom}, the per-turn gain of identical
random and popularity question orders run \emph{under each recommender},
which tests whether the learned instrument suppresses or preserves the
effect the testbed exists to measure.

\begin{table}[t]
\centering
\caption{Learned instrument vs.\ matrix-factorization references, identical
$n{=}200$ held-out users and reveal splits ($d{=}24$ factors; held-out 40\%
of each user's ratings; full-slate ranking of held-out liked items against
non-interacted movies). Explicit rating-MF is a poor top-$N$ ranker by
construction (no negative signal); WRMF (implicit ALS) is therefore the fair
ranking reference. The instrument beats both classical recommenders on every
ranking metric and on calibration, and reaches the highest absolute quality
under elicitation --- so it is neither weak nor masking the measured effect.
Elicitation-AUC = mean NDCG@10 over a 15-turn dialogue under that recommender.}
\label{tab:mf_reference}
\small
\begin{tabular}{lcccc cc}
\toprule
 & \multicolumn{3}{c}{Recommendation (full reveal)} & Recon. & \multicolumn{2}{c}{Elicit-AUC} \\
Recommender & NDCG@10 & Recall@10 & MRR & ROC-AUC & random & pop. \\
\midrule
Popularity-rank (non-pers.) & .175 & .264 & .155 & .596 & --- & --- \\
Explicit MF (rating recon.) & .010 & .026 & .019 & .619 & .010 & .018 \\
WRMF (implicit, ranker)     & .198 & .310 & .170 & .672 & .030 & .113 \\
\textbf{Learned instrument (ours)} & \textbf{.351} & \textbf{.471} & \textbf{.303} & \textbf{.681} & \textbf{.199} & \textbf{.226} \\
\bottomrule
\end{tabular}
\end{table}

\subsection{Metrics}
\label{sec:metrics}

\paragraph{Primary.} The per-turn accuracy curve
$(\mathrm{Acc}_0, \ldots, \mathrm{Acc}_T)$ (Eq.~\ref{eq:acc}), averaged
over users with bootstrap 95\% confidence intervals; \emph{final accuracy}
$\mathrm{Acc}_T$; and \emph{AUAC}, the mean of the curve over turns, which
rewards front-loaded information gain --- two policies with equal final
accuracy differ in AUAC if one achieves its gains earlier. BCE
trajectories are reported alongside as the smooth counterpart of accuracy.

\paragraph{Behavioural diagnostics.} Four quantities decompose \emph{why}
a policy scores as it does:
\begin{description}
  \item[Hit rate.] The fraction of questions answered other than
  \textsf{unknown}. Separates ``asks about things this user can answer''
  from ``asks about informative things''; the popularity baseline
  maximises it by construction.
  \item[Redundancy rate.] The fraction of questions whose answer was
  already inferable: the instrument's pre-question prediction for the
  asked entity exceeded $0.8$ (or fell below $0.2$) and the answer agreed.
  A redundant question is locally safe but strategically wasted.
  \item[Question-type mix.] The distribution of question targets over
  \{genre, actor, director, movie\} as a function of turn number ---
  whether a policy opens broad and refines, or interrogates items
  immediately.
  \item[Coverage.] The entropy of the policy's question distribution
  across users at each turn: a fixed-sequence policy has zero coverage
  entropy at every turn (every user gets the same question), a fully
  user-adaptive policy a high one. This is the metric under which the
  fixed-policy collapse of Section~\ref{sec:casestudy} is visible at a
  glance.
\end{description}

\paragraph{LLM accounting.} For LLM policies we report parse-failure rates
and API cost. Failed parses fall back to a random selection and are
counted against the policy, so formatting unreliability degrades LLM
scores exactly as it would degrade a deployed system.

\paragraph{Statistical reporting.} All policy contrasts are reported as
paired differences over the identical user set with bootstrap CIs;
the instrument's measured per-reveal gain and between-user dispersion
(T4) accompany every table so that readers can verify the contrasts are
resolvable at the sample size used (Section~\ref{sec:snr_theory}).

\section{Reference Policies}
\label{sec:policies}

All policies select from the same slate and observe only the dialogue
record --- never the user profile or population statistics derived from
evaluation users.

\subsection{Non-Learned Baselines}

\textbf{Random} selects uniformly among unasked entities, lower-bounding
the value of any strategy. \textbf{Popularity} asks entities in descending
population rating frequency (computed on training users only): a maximally
answerable, strategy-free order that isolates the value of hit rate alone.
\textbf{Greedy information gain} performs one-step lookahead under the
instrument: for each candidate entity, the expected reduction in the total
binary entropy of the instrument's movie predictions, with the expectation
taken over the candidate's answer --- the instrument's own current
prediction for the entity supplies $p(\textsf{liked})$, and the
population rated-frequency supplies $p(\textsf{answerable})$:
\begin{equation}
\mathrm{EIG}(e) = p_{\mathrm{rated}}(e)\,\Big[\, p_\ell(e)\,
\big(H_t - H_t^{e \mapsto \textsf{liked}}\big) + (1 - p_\ell(e))\,
\big(H_t - H_t^{e \mapsto \textsf{disliked}}\big) \Big].
\end{equation}
This transplants the classical acquisition principle
\citep{viappiani2009regret, christakopoulou2016towards} onto the testbed
instrument and is the benchmark any learned or prompted policy must beat to
claim strategic value.

\subsection{Prompt-Only LLM Questioners}

Three styles span the prompting sophistication a practitioner might
reasonably deploy, all menu-constrained to the slate and parsed with a
lenient matcher: \textbf{vanilla} (choose the next entity to ask about);
\textbf{strategist} (a brief chain-of-thought prompt articulating
elicitation principles: broad attributes first, follow up on what was
learned, avoid unanswerable and redundant questions); and \textbf{GATE}
(framing from \citet{li2023eliciting}: conduct an expert interview, probe
edge cases and contrasts between taste hypotheses). Decoding is
deterministic (temperature 0). \todo{models: gpt-4o-mini primary; consider
one frontier model for the camera-ready}

\subsection{Learned Bot-Play Policy}

We retrain the bot-play REINFORCE method of \citet{makarova2024learning} on
the testbed: a feedforward policy over the revealed-state encoding, trained
on training-partition users with per-turn reward equal to the reduction in
the instrument's BCE loss, $\epsilon$-greedy exploration, and a running
per-turn baseline. This is a faithful transplantation of the published
method onto the shared instrument, making the comparison with prompted and
information-theoretic policies like-for-like for the first time.

\section{Experimental Setup}
\label{sec:setup}

\todo{finalise numbers before results runs}
Evaluation uses $n = 200$ testbed users (sampled with seed 42 from the
reserved partition), a turn budget of $T = 15$, and identical user sets and
seeds for every policy. The instrument is \textsc{OneHot+RS}
(Section~\ref{sec:reveal_subset}); instrument ablations re-run the full
benchmark under the rejected \textsc{OneHot} instrument to quantify how
instrument failure distorts policy rankings. LLM policies use
\todo{gpt-4o-mini, version-pinned}; total API cost for the benchmark is
\results{\$X}.

\section{Results}
\label{sec:results}

\subsection{Main Leaderboard (Slate 1)}

Table~\ref{tab:leaderboard} reports the primary comparison: twelve
policies spanning every elicitation family in the related work, evaluated
on 200 held-out users at a 15-turn budget under the accepted instrument.
\emph{Answer rate} is the fraction of questions the simulated user could
answer (renamed from ``hit rate'' to avoid collision with Hit@$k$);
\emph{adaptive} reports the behavioural audit verdict (Section~\ref{sec:metrics}):
whether the policy's question sequence branches on the user's answers.

\begin{table}[t]
\centering
\caption{Slate-1 leaderboard, 200 held-out users, $T=15$.
AUAC@$k$ = mean accuracy over the first $k$ turns. Sorted by AUAC@15.
\todo{PPO and DQN rows pending the extended convergence runs; paired
bootstrap CIs for new rows; llm\_gate tail re-run}}
\label{tab:leaderboard}
\small
\begin{tabular}{lcccccc}
\toprule
Policy & AUAC@5 & AUAC@15 & Final & NDCG@10 & Answer rate & Adaptive \\
\midrule
Bot-play RL \citep{makarova2024learning} & .724 & \textbf{.752} & .779 & .690 & 71\% & no \\
PPO & .727 & .751 & .778 & .686 & 66\% & no$^{\dagger}$ \\
SCPR-entropy \citep{lei2020interactive} & .727 & .751 & \textbf{.785} & .625 & 76\% & yes \\
Greedy info-gain & .726 & .747 & .776 & .657 & 69\% & yes \\
Bot-play RL v2 & .725 & .742 & .764 & .637 & 81\% & no \\
DQN \citep{deng2021unified} & \textbf{.727} & .740 & .753 & .613 & 44\% & yes \\
Popularity & .725 & .739 & .757 & .608 & 88\% & no \\
Random & .716 & .735 & .757 & .616 & 37\% & --- \\
LLM strategist (gpt-4o-mini) & .709 & .730 & .763 & \textbf{.708} & 64\% & yes \\
LLM vanilla (gpt-4o-mini) & .707 & .728 & .752 & .635 & 40\% & yes \\
Thompson \citep{austin2024bayesian} & .712 & .727 & .749 & .693 & 49\% & yes \\
LLM GATE-style (gpt-4o-mini) & .707 & .726 & .750 & .635 & 40\% & yes \\
\bottomrule
\end{tabular}
\end{table}

Five findings. \textbf{(1) Prompt-only LLMs do not ask strategically:}
all three styles rank below greedy information gain at every turn budget,
and two rank below \emph{random}; their first two questions gain almost
no accuracy ($+0.002$ by turn 2 versus $+0.022$ for the top tier). This
quantifies, for the first time under controlled conditions, the
elicitation weakness attributed to LLMs in the literature
(Section~\ref{sec:related}). \textbf{(2) The learned bot-play policy is
top tier} --- but the behavioural audit shows it (and PPO) achieve this
with \emph{static} question sequences, while every audited-adaptive
policy scores at or below them. On a popularity-saturated slate, the
best questions are nearly the same for everyone; the audit, not the
score, reveals this. \textbf{(3) PPO chooses staticity:} the same
algorithm that provably learns adaptive routing on the synthetic world
(Section~\ref{sec:synthetic_results}) converges to a playlist here ---
evidence that staticity is a property of the slate, not a failure of
learning. \textbf{(4) Metric choice reshuffles the ranking:} on NDCG@10
the LLM strategist is \emph{best} (.708) and Thompson second (.693) ---
policies that pursue likes optimise ranking, policies that pursue
uncertainty optimise classification. Elicitation benchmarks must report
both. \textbf{(5) Short budgets invert the order:} at AUAC@3, popularity
leads --- answerability dominates the opening turns; strategy pays only
from turn $\sim$5. \results{paired bootstrap table; per-policy CIs}

\subsection{Elicitation Efficiency --- The Right Effect Size}
\label{sec:efficiency}

Final-turn accuracy understates elicitation value, because at a long
budget every competent policy converges toward the instrument's
information ceiling and the endpoint gaps compress. The quantity that
matters operationally is \emph{efficiency}: how few questions a policy
needs to reach a given recommendation quality, since each question spends
scarce user patience. We therefore report, for every policy, the number
of turns to match the accuracy that \emph{random} questioning reaches at
the full 15-turn budget, and the implied speed-up
(Table~\ref{tab:efficiency}); the per-turn curves are
Figures~\ref{fig:eff1}--\ref{fig:eff2}.

\begin{table}[t]
\centering
\caption{Elicitation efficiency. ``turns@R'' = turns to reach random's
15-turn accuracy; speed-up $=15/\text{turns@R}$; ``—'' = never reaches it
within 15 turns. Measured this way, strategy yields a 2.1--2.5$\times$
question-efficiency gain on both slates.}
\label{tab:efficiency}
\small
\begin{tabular}{lccc|ccc}
\toprule
 & \multicolumn{3}{c|}{\textbf{Slate 1} (prior .700)}
 & \multicolumn{3}{c}{\textbf{Slate 2} (prior .638)} \\
Policy & AUAC@5 & turns@R & speed-up & AUAC@5 & turns@R & speed-up \\
\midrule
Random            & .716 & 15 & 1.0$\times$ & .661 & 15 & 1.0$\times$ \\
Popularity        & .725 & 14 & 1.1$\times$ & .678 & — & — \\
Greedy info-gain  & .726 & 10 & 1.5$\times$ & .687 & 11 & 1.4$\times$ \\
SCPR-entropy      & .727 & 9  & 1.7$\times$ & .683 & 7  & 2.1$\times$ \\
Bot-play \citep{makarova2024learning} & .724 & 7 & \textbf{2.1$\times$} & .686 & 6 & \textbf{2.5$\times$} \\
PPO+ans           & .725 & 7  & \textbf{2.1$\times$} & .685 & 6  & \textbf{2.5$\times$} \\
\bottomrule
\end{tabular}
\end{table}

Three points follow. First, \textbf{the effect is real and sizeable when
measured correctly}: the learned and uncertainty-driven policies reach in
6--7 questions what random reaches in 15 --- a 2.1--2.5$\times$
efficiency gain --- whereas the same comparison on final accuracy alone
spans under two points and reads as negligible. Second,
\textbf{static non-adaptive ordering does not get this for free}:
popularity never reaches random's 15-turn accuracy on slate 2 (it
saturates on broadly-answerable but low-information questions), so the
gain is attributable to \emph{what} is asked, not merely to asking
answerable things. Third, \textbf{the gap widens at short budgets}: at
turn 5 the strategy--random accuracy gap is roughly double its turn-15
value, and real elicitation dialogues are short (3--6 turns), precisely
the regime where efficiency matters most. We therefore treat
turns-to-quality, not endpoint accuracy, as the primary effect size
throughout.

\begin{figure}[t]
\centering
\includegraphics[width=0.78\textwidth]{figures/fig_efficiency_slate1.pdf}
\caption{Per-turn accuracy, slate 1. Dotted line: random's 15-turn level;
strategy reaches it in $\sim$7 turns. \results{auto-regenerated by
\texttt{efficiency\_analysis.py}}}
\label{fig:eff1}
\end{figure}

\begin{figure}[t]
\centering
\includegraphics[width=0.78\textwidth]{figures/fig_efficiency_slate2.pdf}
\caption{Per-turn accuracy, slate 2 (mid-popularity). The efficiency gap
is larger than slate 1: strategy reaches random's endpoint in $\sim$6
turns; popularity never does.}
\label{fig:eff2}
\end{figure}

\subsection{The Synthetic Indicator World}
\label{sec:synthetic_results}

\begin{table}[t]
\centering
\caption{Synthetic indicator world (300 users, $T=5$; instrument gates:
oracle-reveal ceiling 0.904, answerability separation $+0.71$). Every
baseline class fails to capture the proven adaptive headroom for a
different, diagnosable reason; PPO captures 44\% of it and exceeds the
hand-built oracle on AUAC by avoiding the indicator's zero-immediate-gain
turn.}
\label{tab:synthetic}
\small
\begin{tabular}{lccccl}
\toprule
Policy & AUAC & Final & Branches & Asks indicator & Diagnosis \\
\midrule
Oracle-adaptive & .623 & \textbf{.804} & yes & 100\% & analytic ceiling \\
PPO & \textbf{.628} & .742 & \textbf{yes} & 0\% & learned, adaptive \\
DQN & .623 & .738 & \textbf{yes} & 0\% & learned, adaptive \\
SCPR-entropy & .609 & .710 & yes & 0\% & heuristic, adaptive \\
Static oracle & .591 & .694 & no & 0\% & best non-adaptive \\
REINFORCE-v3 & .591 & .662 & \textbf{no} & 0\% & learned, \textbf{static} \\
Greedy + answerability & .584 & .664 & yes & 0\% & myopia-bound \\
Greedy info-gain & .582 & .660 & yes & 0\% & myopia-bound \\
Random & .583 & .657 & --- & 6\% & --- \\
REINFORCE-v2 & .574 & .619 & no & 0\% & learned, static \\
Popularity & .554 & .618 & no & 100\% & static; wastes indicator \\
\bottomrule
\end{tabular}
\end{table}

Three structural results. First, the \emph{answerability discovery}: with
a single-head instrument exposing only $P(\textsf{liked})$, every policy
failed ($\leq$0.59 AUAC against a verified 0.90 ceiling) because the
world's adaptive value flows through answerability routing, which no
policy could observe; the dual head repairs expressibility.
Second, \emph{myopia is structural}: the indicator has zero one-step
information gain --- its value is purely future routing --- so no amount
of belief quality rescues one-step lookahead. Third, \emph{the optimiser
is the boundary, not the framework}: the adaptive/static divide in the
table is near-perfect, and the only learned policy on the static side of
it is REINFORCE (the bot-play lineage's optimiser, including
\citealp{das2017learning, makarova2024learning}), which saturates exactly
at the optimal \emph{static} policy. Given the same state, reward, and
budget, \emph{both} modern optimisers cross the line --- PPO
(advantage-based) and DQN (value-based with replay) each learn to branch
and approach the analytic oracle, PPO discovering a routing trick
(informative probes as implicit indicators) superior to the hand-designed
oracle's opening move. The implication for the bot-play programme is
direct: its ceiling was the high-variance REINFORCE update, not the
framework; a modern optimiser unlocks adaptive elicitation. (Popularity is
the instructive failure: it asks the always-answerable indicator first
100\% of the time, yet scores worst, because a static policy cannot act on
the answer it receives.)

\subsection{Remaining artefacts}

\results{Generated by \texttt{make\_figures.py} from episode logs; thesis
version carries the full set:}

\subsubsection*{Main comparison}
\begin{itemize}
  \item \textbf{Table 5 --- Main leaderboard.} Final accuracy (CI), AUAC
  (CI), hit rate, redundancy rate, per policy; paired differences vs
  greedy information gain. \results{}
  \item \textbf{Figure 1 --- Accuracy curves.} Mean per-turn accuracy with
  CI bands, all eight policies, $T = 15$. \results{}
  \item \textbf{Figure 2 --- BCE curves.} The smooth counterpart of
  Figure~1; separates policies whose accuracy ties mask confidence
  differences. \results{}
\end{itemize}

\subsubsection*{Behavioural analysis: where strategies differ}
\begin{itemize}
  \item \textbf{Figure 3 --- Hit rate by turn.} Which policies waste turns
  on unanswerable questions, and when. \results{}
  \item \textbf{Figure 4 --- Question-type mix.} Stacked per-turn
  composition (genre/actor/director/movie) per policy: opening strategy
  and refinement behaviour at a glance. \results{}
  \item \textbf{Figure 5 --- Question--turn heatmap.} For each policy, the
  20 most-asked entities $\times$ turn position: fixed sequences appear as
  vertical bars, adaptive policies as diffuse mass. The direct visual
  test of the fixed-policy pathology of Section~\ref{sec:casestudy}.
  \results{}
  \item \textbf{Figure 6 --- Coverage entropy by turn.} Quantifies
  Figure~5: question-distribution entropy across users per turn.
  \results{}
  \item \textbf{Figure 7 --- Redundancy by turn.} When policies start
  asking questions whose answers they should already infer. \results{}
  \item \textbf{Table 6 --- LLM behaviour.} Parse failures, cost, and the
  behavioural metrics above contrasted across the three prompt styles:
  does prompting sophistication buy strategy or only fluency? \results{}
\end{itemize}

\subsubsection*{Validity and robustness}
\begin{itemize}
  \item \textbf{Figure 8 --- Instrument ablation.} The full benchmark
  re-run under the rejected full-profile instrument: policy separation
  compresses and ranks distort, the paper's measurement-validity argument
  made visual. \results{}
  \item \textbf{Figure 9 --- Per-user outcome distributions.} Violin plots
  of final accuracy per policy: do strategies differ in mean only, or in
  tail behaviour (users they fail entirely)? \results{}
  \item \textbf{Table 7 --- Turn-budget ablation.} AUAC and final accuracy
  at $T \in \{5, 10, 15\}$: which policies front-load. \results{}
  \item \textbf{Figure 10 --- Power analysis.} Detectable AUAC difference
  vs number of evaluation users at the instrument's measured dispersion;
  justifies $n = 200$ and tells future users of the testbed what their
  $n$ buys. \results{}
  \item \textbf{Figure 11 --- Acceptance-test evolution.} The four
  instrument generations of Table~\ref{tab:audit} as grouped bars; the
  measurement-first narrative in one figure. (Data already complete for
  three generations.) \results{}
  \item \textbf{Table 8 --- Seed robustness.} Main leaderboard re-run
  under three user-sample seeds (free policies; LLM policies one seed,
  cost). \results{}
\end{itemize}

\subsection{Analysis Questions the Results Must Answer}

\begin{enumerate}
  \item Does any LLM style beat greedy information gain on AUAC? By how
  much, and does the answer change with turn budget?
  \item Where do LLM questioners lose: hit rate (asking unanswerable
  things), redundancy (asking inferable things), or breadth (failing to
  cover the taste space)?
  \item Does the learned bot-play policy recover greedy's performance
  without access to explicit lookahead, and does it exploit population
  structure greedy cannot see?
  \item How much of the policy ranking survives under the broken
  instrument? (If little: quantified evidence that instrument acceptance
  must precede benchmarking.)
\end{enumerate}

\section{Discussion}
\label{sec:discussion}

\results{To be written from results. Scaffolding for the two principal
outcome branches:}

\paragraph{If greedy information gain dominates LLM questioning.} The
field's informal claim becomes a measured gap: prompt-only LLMs, however
fluent, do not implement effective elicitation strategy, and the gap
justifies the learned-policy research agenda --- with the testbed supplying
the training signal and the evaluation. The interesting quantity is the
decomposition (hit rate vs redundancy vs breadth), which tells
\emph{prompting} research what to fix.

\paragraph{If LLM questioning matches or beats greedy.} The motivating
claim of several 2024--2026 systems needs revision: the binding constraint
in LLM-era CRS elicitation is not strategy selection but the surrounding
measurement and integration. The value of the testbed inverts but does not
vanish: it becomes the evidence that cheap prompting suffices, and the
learned-policy agenda should target the regimes where it still fails
(long-horizon, niche tastes, low hit-rate slates).

\section{Limitations}
\label{sec:limitations}

The testbed measures strategy under deliberate abstractions. The simulator
answers labels, not language: linguistic skill --- phrasing questions users
understand and enjoy, parsing free-form answers --- is excluded by design
and must be evaluated by other means \citep{kostric2025should}; CASPER
licenses claims about \emph{what} to ask, not \emph{how} to ask it. The
slate is closed and movie-domain; open-vocabulary elicitation, the subject
of our companion work, requires extending the instrument to unseen
concepts. Attribute ground truth is derived from movie ratings rather than
stated directly, inheriting that approximation. Binarised ratings discard
preference intensity. Finally, simulated users drawn from MovieLens
profiles represent the population of historical raters, not of
cold-start users in deployment.

\section{Conclusion}
\label{sec:conclusion}

We introduced CASPER, a controlled testbed isolating strategic preference
elicitation in conversational recommendation, built on an audited
measurement instrument, a verifiable simulator, and a shared action slate.
Our instrument audit shows that measurement validity is not free: plausible
configurations fail polarity or reveal-regime requirements in ways that
would silently corrupt any policy comparison run on top of them, and a
simple reveal-subset augmentation repairs the failure. The benchmark
provides the first like-for-like comparison spanning random,
information-theoretic, prompted-LLM, and reinforcement-learned elicitation.
\results{one-sentence headline finding}. The released testbed turns a
contested capability claim into a measurable quantity, and provides the
fixed training and evaluation signal on which our companion work on
continuous-action elicitation policies builds.

\paragraph{Acknowledgements.} \todo{}

\bibliographystyle{agsm}
\bibliography{references}

\appendix

\section{LLM Questioner Prompts}
\label{app:prompts}

The three prompt styles of Section~\ref{sec:policies}, verbatim. All use
temperature 0 and a 150-token completion limit; the menu of remaining
entities and the dialogue transcript are appended to each request.

\subsection*{A.1 Vanilla}
\begin{verbatim}
You are a movie recommendation assistant interviewing a user to
learn their taste. At each step you choose ONE entity to ask about
(whether they like it), from a fixed menu. Reply with the exact
entity name from the menu, and nothing else.
\end{verbatim}

\subsection*{A.2 Strategist}
\begin{verbatim}
You are an expert preference-elicitation strategist for a movie
recommender. Your goal: in as few questions as possible, learn
enough about this user's taste to predict their ratings for 100
popular movies. At each step you choose ONE entity from a fixed
menu.

Think strategically before answering:
- Broad attributes (genres) split the space fastest early on.
- Follow up on what you have learned: refine within liked areas,
  confirm dislikes.
- Questions the user can't answer ('unknown') are wasted turns;
  popular entities are safer.
- Avoid redundant questions whose answer you can already infer.

First think step by step in at most 50 words, then output the
chosen entity name on the final line in the exact form:
CHOICE: <entity name>
\end{verbatim}

\subsection*{A.3 GATE-style}
\begin{verbatim}
You are conducting a free-form preference interview, but your
questions must be realised by selecting ONE entity from a fixed
menu (the system asks 'do you like <entity>?'). Run the interview
the way an expert interviewer would: open with broad probes, test
edge cases of what you believe, and probe contrasts (if they like
X, check a nearby Y that would distinguish two hypotheses about
their taste). Output the chosen entity name on the final line as:
CHOICE: <entity name>
\end{verbatim}

\section{Instrument Training and Audit Detail}
\label{app:instrument}

\subsection*{B.1 Training configuration}
Architecture: item embedding (dimension $\lfloor N/2 \rfloor = 93$)
concatenated with the rating one-hot, single-layer LSTM (hidden 93),
dense bottleneck, single-head attention, linear head over the slate;
193{,}535 parameters. Optimiser Adam, learning rate $10^{-3}$, batch 32,
masked BCE over rated targets at every sequence position, early stopping
on validation loss (patience 10). Reveal-subset augmentation: per sample,
$k \sim \exp(\mathcal{U}(0, \ln(n_{\mathrm{rated}}+1)))$ revealed items;
with probability 0.15 the candidate reveal set is restricted to
attributes. Training population: 12{,}000 users of the 80\% partition
(of 64{,}639 eligible); model selection on 1{,}500 held-out users.
Wall-clock on a 4-core laptop CPU: $\approx$45--60 s per epoch,
early-stopped at \results{epoch} ($\approx$\results{} minutes total).
\results{training curve figure}

\subsection*{B.2 Franchise flip detail}
\results{per-franchise table: anchor, related movies, rank movement under
liked vs disliked anchor, for the accepted instrument --- the v1 audit
measured 8 franchises with slate support (Star Wars, Matrix, Toy Story,
Lord of the Rings, Terminator, Back to the Future, Jurassic Park,
Godfather, et al.)}

\subsection*{B.3 Attribute label statistics}
\results{coverage: fraction of users carrying a non-unknown label per
attribute under absolute vs taste-relative derivation; the taste-relative
scheme abstains where the absolute scheme fabricated severity-driven
labels}

\section{Reproducibility}
\label{app:repro}

All experiments run on a single laptop (Intel i7-7820HQ, 16~GB RAM, CPU
PyTorch 2.2). Data: MovieLens-25M \citep{harper2015movielens}; cast and
crew from TMDB. User partitions derive from one global shuffle (seed 42):
80\% training / first 1{,}500 of the remainder for instrument model
selection / final $\sim$11{,}400 reserved for evaluation; evaluation
samples and bootstrap resampling use fixed seeds recorded in the released
configuration. LLM policies: \todo{model and version pin}, temperature 0;
total benchmark API cost \results{\$X}. The released repository contains
the instrument checkpoints (all four generations), acceptance suite,
simulator, policies, benchmark runner, figure generation, and the episode
logs underlying every figure. \todo{repository URL}

\section{Entity Slate}
\label{app:slate}

The 187-item slate: the 100 most-rated MovieLens movies; 17 genres
(Action, Adventure, Animation, Children, Comedy, Crime, Documentary,
Drama, Fantasy, Film-Noir, Horror, IMAX, Musical, Mystery, Romance,
Sci-Fi, Thriller); the 50 most-frequent lead actors and 20 directors of
those movies. \results{full listing or pointer to repository file}

\end{document}
